%-----------------------------------------------------------------------------------------------------------------------------------------------%
%	The MIT License (MIT)
%
%	Copyright (c) 2015 Jan Küster
%
%	Permission is hereby granted, free of charge, to any person obtaining a copy
%	of this software and associated documentation files (the "Software"), to deal
%	in the Software without restriction, including without limitation the rights
%	to use, copy, modify, merge, publish, distribute, sublicense, and/or sell
%	copies of the Software, and to permit persons to whom the Software is
%	furnished to do so, subject to the following conditions:
%	
%	THE SOFTWARE IS PROVIDED "AS IS", WITHOUT WARRANTY OF ANY KIND, EXPRESS OR
%	IMPLIED, INCLUDING BUT NOT LIMITED TO THE WARRANTIES OF MERCHANTABILITY,
%	FITNESS FOR A PARTICULAR PURPOSE AND NONINFRINGEMENT. IN NO EVENT SHALL THE
%	AUTHORS OR COPYRIGHT HOLDERS BE LIABLE FOR ANY CLAIM, DAMAGES OR OTHER
%	LIABILITY, WHETHER IN AN ACTION OF CONTRACT, TORT OR OTHERWISE, ARISING FROM,
%	OUT OF OR IN CONNECTION WITH THE SOFTWARE OR THE USE OR OTHER DEALINGS IN
%	THE SOFTWARE.
%	
%
%-----------------------------------------------------------------------------------------------------------------------------------------------%


%============================================================================%
%
%	DOCUMENT DEFINITION
%
%============================================================================%

%we use article class because we want to fully customize the page and dont use a cv template
\documentclass[10pt,A4]{article}	


%----------------------------------------------------------------------------------------
%	ENCODING
%----------------------------------------------------------------------------------------

%we use utf8 since we want to build from any machine
\usepackage[utf8]{inputenc}		

%----------------------------------------------------------------------------------------
%	LOGIC
%----------------------------------------------------------------------------------------

% provides \isempty test
\usepackage{xifthen}

%----------------------------------------------------------------------------------------
%	FONT
%----------------------------------------------------------------------------------------

% some tex-live fonts - choose your own

%\usepackage[defaultsans]{droidsans}
%\usepackage[default]{comfortaa}
%\usepackage{cmbright}
\usepackage[default]{raleway}
%\usepackage{fetamont}
%\usepackage[default]{gillius}
%\usepackage[light,math]{iwona}
%\usepackage[thin]{roboto} 

% set font default
\renewcommand*\familydefault{\sfdefault} 	
\usepackage[T1]{fontenc}

% more font size definitions
\usepackage{moresize}		


%----------------------------------------------------------------------------------------
%	PAGE LAYOUT  DEFINITIONS
%----------------------------------------------------------------------------------------

%debug page outer frames
%\usepackage{showframe}			


%define page styles using geometry
\usepackage[a4paper]{geometry}		

% for example, change the margins to 2 inches all round
\geometry{top=1.75cm, bottom=-.6cm, left=1.5cm, right=1.5cm} 	

%use customized header
\usepackage{fancyhdr}				
\pagestyle{fancy}

%less space between header and content
\setlength{\headheight}{-5pt}		


%customize entries left, center and right
\lhead{}
\chead{ \small{Andreas R. Spyridakis  $\cdot$ Aluno de ciência da computação $\cdot$  Brasília, DF  $\cdot$  \textcolor{sectcol}{\textbf{andreas.oxyz@proton.me}}  $\cdot$ +55 61 *****-7578}}
\rhead{}


%indentation is zero
\setlength{\parindent}{0mm}

%----------------------------------------------------------------------------------------
%	TABLE /ARRAY DEFINITIONS
%---------------------------------------------------------------------------------------- 

%for layouting tables
\usepackage{multicol}			
\usepackage{multirow}

%extended aligning of tabular cells
\usepackage{array}

\newcolumntype{x}[1]{%
>{\raggedleft\hspace{0pt}}p{#1}}%


%----------------------------------------------------------------------------------------
%	GRAPHICS DEFINITIONS
%---------------------------------------------------------------------------------------- 

%for header image
\usepackage{graphicx}

%for floating figures
\usepackage{wrapfig}
\usepackage{float}
%\floatstyle{boxed} 
%\restylefloat{figure}

%for drawing graphics		
\usepackage{tikz}				
\usetikzlibrary{shapes, backgrounds,mindmap, trees}


%----------------------------------------------------------------------------------------
%	Color DEFINITIONS
%---------------------------------------------------------------------------------------- 

\usepackage{color}

%accent color
\definecolor{sectcol}{RGB}{255,150,0}

%dark background color
\definecolor{bgcol}{RGB}{110,110,110}

%light background / accent color
\definecolor{softcol}{RGB}{225,225,225}


%----------------------------------------------------------------------------------------
%	HYPERLINKS DEFINITIONS
%---------------------------------------------------------------------------------------- 
\usepackage[hidelinks]{hyperref}
\hypersetup{
    colorlinks=false,
    linkcolor=blue,
    filecolor=magenta,      
    urlcolor=blue,
    pdftitle={Overleaf Example},
    pdfpagemode=FullScreen,
    }

\urlstyle{same}

%============================================================================%
%
%
%	DEFINITIONS
%
%
%============================================================================%

%----------------------------------------------------------------------------------------
% 	HEADER
%----------------------------------------------------------------------------------------

% remove top header line
\renewcommand{\headrulewidth}{0pt} 

%remove botttom header line
\renewcommand{\footrulewidth}{0pt}	  	

%remove pagenum
\renewcommand{\thepage}{}	

%remove section num		
\renewcommand{\thesection}{}			

%----------------------------------------------------------------------------------------
% 	ARROW GRAPHICS in Tikz
%----------------------------------------------------------------------------------------

% a six pointed arrow poiting to the left
\newcommand{\tzlarrow}{(0,0) -- (0.2,0) -- (0.3,0.2) -- (0.2,0.4) -- (0,0.4) -- (0.1,0.2) -- cycle;}	

% include the left arrow into a tikz picture
% param1: fill color
%
\newcommand{\larrow}[1]
{\begin{tikzpicture}[scale=0.58]
	 \filldraw[fill=#1!100,draw=#1!100!black]  \tzlarrow
 \end{tikzpicture}
}

% a six pointed arrow poiting to the right
\newcommand{\tzrarrow}{ (0,0.2) -- (0.1,0) -- (0.3,0) -- (0.2,0.2) -- (0.3,0.4) -- (0.1,0.4) -- cycle;}

% include the right arrow into a tikz picture
% param1: fill color
%
\newcommand{\rarrow}
{\begin{tikzpicture}[scale=0.7]
	\filldraw[fill=sectcol!100,draw=sectcol!100!black] \tzrarrow
 \end{tikzpicture}
}



%----------------------------------------------------------------------------------------
%	custom sections
%----------------------------------------------------------------------------------------

% create a coloured box with arrow and title as cv section headline
% param 1: section title
%
\newcommand{\cvsection}[1]
{
\colorbox{sectcol}{\mystrut \makebox[1\linewidth][l]{
\larrow{bgcol} \hspace{-8pt} \larrow{bgcol} \hspace{-8pt} \larrow{bgcol} \textcolor{white}{\textbf{#1}}\hspace{4pt}
}}\\
}

%create a coloured arrow with title as cv meta section section
% param 1: meta section title
%
\newcommand{\metasection}[2]
{
\begin{tabular*}{1\textwidth}{p{2.4cm} p{11cm}}
\larrow{bgcol}	\normalsize{\textcolor{sectcol}{#1}}&#2\\[12pt]
\end{tabular*}
}

%----------------------------------------------------------------------------------------
%	 CV EVENT
%----------------------------------------------------------------------------------------

% creates a stretched box as cv entry headline followed by two paragraphs about 
% the work you did
% param 1:	event time i.e. 2014 or 2011-2014 etc.
% param 2:	event name (what did you do?)
% param 3:	institution (where did you work / study)
% param 4:	what was your position
% param 5:	some words about your contributions
%
\newcommand{\cvevent}[5]
{
\vspace{8pt}
	\begin{tabular*}{1\textwidth}{p{2.3cm}  p{10.8cm} x{3.9cm}}
 \textcolor{bgcol}{#1}& \textbf{#2} & \vspace{2.5pt}\textcolor{sectcol}{#3}

	\end{tabular*}
\vspace{-12pt}
\textcolor{softcol}{\hrule}
\vspace{6pt}
	\begin{tabular*}{1\textwidth}{p{2.3cm} p{14.4cm}}
&		 \larrow{bgcol}  #4\\[3pt]
&		 \larrow{bgcol}  #5\\[6pt]
	\end{tabular*}

}

% creates a stretched box as 
\newcommand{\cveventmeta}[2]
{
	\mbox{\mystrut \hspace{87pt}\textit{#1}}\\
	#2
}

%----------------------------------------------------------------------------------------
% CUSTOM STRUT FOR EMPTY BOXES
%----------------------------------------- -----------------------------------------------
\newcommand{\mystrut}{\rule[-.3\baselineskip]{0pt}{\baselineskip}}

%----------------------------------------------------------------------------------------
% CUSTOM LOREM IPSUM
%----------------------------------------------------------------------------------------
\newcommand{\lorem}
{Lorem ipsum dolor sit amet, consectetur adipiscing elit. Donec a diam lectus.}



%============================================================================%
%
%
%
%	DOCUMENT CONTENT
%
%
%
%============================================================================%
\begin{document}


%use our custom fancy header definitions
\pagestyle{fancy}	


%---------------------------------------------------------------------------------------
%	TITLE HEADLINE
%----------------------------------------------------------------------------------------
\vspace{-20.55pt}

% use this for multiple words like working titles etc.
%\hspace{-0.25\linewidth}\colorbox{bgcol}{\makebox[1.5\linewidth][c]{\hspace{46pt}\HUGE{\textcolor{white}{\textsc{Jan Küster}} } \textcolor{sectcol}{\rule[-1mm]{1mm}{0.9cm}} \parbox[b]{5cm}{   \large{ \textcolor{white}{{IT Consultant}}}\\
% \large{ \textcolor{white}{{Resume}}}}
%}}

% use this for single words, e.g. CV or RESUME etc.
\hspace{-0.25\linewidth}\colorbox{bgcol}{\makebox[1.5\linewidth][c]{\huge{\textcolor{white}{\textsc{Andreas Ramalho Spyridakis}} } \textcolor{sectcol}{\rule[-1mm]{1mm}{0.9cm}} \huge{\textcolor{white}{\textsc{Curriculo Vitae}} } }}


%----------------------------------------------------------------------------------------
%	HEADER IMAGE
%----------------------------------------------------------------------------------------

\begin{figure}[H]
\begin{flushright}
	%\includegraphics[trim= 320 130 460 210,clip,width=0.2\linewidth]{myfoto.jpg}	%trimming relative to image size!
    \includegraphics[trim= 10 25 0 0,clip,width=0.2\linewidth]{myfoto.jpg}	%trimming relative to image size!
\end{flushright}
\end{figure}

%---------------------------------------------------------------------------------------
%	QR CODE (optional)
%----------------------------------------------------------------------------------------
%\vspace{-136pt}
%\hspace{0.75\linewidth}
%\includegraphics[width=103pt]{qrcode}
%\normalsize
%\vspace{88pt}

%---------------------------------------------------------------------------------------
%	META SECTION
%----------------------------------------------------------------------------------------

\vspace{-114pt}

\metasection{Situação:}{Aluno de B.Sc. Ciência da computação}
%\metasection{Status:}{Aluno de B.Sc. Ciência da computação}
\metasection{Campos:}{Suporte de TI, Desenvolvimento de software}
%\metasection{Fields:}{Suporte de TI, Desenvolvimento de software}
\metasection{Tecnologias:}{Python, Javascript, HTML, Markdown, CSS, Latex, Astro, Hugo, Jekyll, Git, Vim, Emacs, Visual Studio Code, EclipseIDE, MongoDB, MySQL/MariaDB}
%\metasection{Tech:}{Python, Javascript, HTML, CSS, MongoDB, MySQL/MariaDB, Git}
\metasection{Linguas:}{Português (nativo), Inglês (avançado)}
%\metasection{Languages:}{Português (native), Inglês (advanced)}
\metasection{Lazer:}{Jogos (simulação e retro), Esportes e preparo físico, Ficção científica}
%\metasection{Hobbies:}{Jogos (simulação e retro), Esportes e preparo físico, Ficção científica}

\vspace{6pt}

%---------------------------------------------------------------------------------------
%	SUMMARAY (optional)
%----------------------------------------------------------------------------------------

%\cvsection{Summary}\\
\cvsection{\large{SUMÁRIO}}\\
Aluno de B.Sc. de Ciência da computação no Instituto de Ensino Superior de Brasília (IESB) com experiência prévia de suporte de TI em órgão da Administração Pública local do Distrito Federal vinculado ao Poder Executivo.\\

Um currículo mais resumido está disponível na seguinte página: \href{https://github.com/andreas0r0s/cv/blob/master/R%C3%A9sum%C3%A9/R%C3%A9sum%C3%A9.pdf}{\underline{\texttt{link de acesso}}}.

%I am a digital media graduate (M.Sc.) with project experience in educational research as well as in the private sector. During my studies I focused on e-assessment software and moved over to b2b software for IBM Notes Domino.

%Currently I develop and evaluate the next generation learning management system with Meteor based on an extensive nursing curriculum for healthcare education.  I also love fitness, martial arts, videogames, news and Sci-Fi series.\\[-2pt]

\vspace{6pt}
%============================================================================%
%
%	CV SECTIONS AND EVENTS (MAIN CONTENT)
%
%============================================================================%

%---------------------------------------------------------------------------------------
%	EXPERIENCE
%----------------------------------------------------------------------------------------
\cvsection{\large{EXPERIÊNCIA}}
%\cvsection{Experience}

%
\cvevent{2022 - 2024}{Estagiário de TI da Administração Regional do Cruzeiro}{Brasília, DF}{Estágio no Núcleo de Informática da RA Cruzeiro de \textit{\textbf{maio de 2022 a maio de 2024}} com suporte técnico e operacional de TI}{
    \textbf{Atividades realizadas}:
    \begin{itemize}
        \item Instalação de redes, microeletrônicos e equipamentos de comunicação
        \item Auxiliar com a \textbf{configuração e manutenção de equipamentos de TI} (impressoras da rede,  microeletrônicos e demais equipamentos da rede local)
        \item Suporte ao usuário
        \item Elaborar relatórios de demandas atendidas e atividades realizadas
        \item Atualizar e documentar informações, dados e fluxos de trabalho da área de TI da Administração (rede, suporte ao usuário, equipamentos de TI e software)
        \item \textbf{Elaboração} do \href{https://cruzeiro.df.gov.br/wp-conteudo/uploads/2024/04/PDTI-RA-CRUZ-\%E2\%80\%93-2023\_2026.pdf}{\textbf{PDTI da Administração do Cruzeiro de 2023-2026}}
        \item \textbf{Levantamento, registro e atualização} do \textbf{inventário de TI} da Administração
        \item \textbf{Troca} de \textbf{100 microcomputadores} obsoletos e defasados por MiniPC's mais atualizados
        \item \textbf{Levantamento de especificações técnicas} e \textbf{criação} de documentos de \textbf{requisição de compra} de \textbf{switches} e de \textbf{nobreaks}
        \item \textbf{Montagem e reparo} de \textbf{equipamentos de TI e cabeamento estruturado} da \textbf{Sala de Estudos} e da \textbf{Sala de Informática do Idoso}
        \item \textbf{Preparo} da \textbf{sala de suporte} para os \textbf{atendimentos do DF-LEGAL} dentro da Administração
        \item \textbf{Montagem e configuração} de \textbf{estação de trabalho} para o \textbf{atendimento} itinerante da \textbf{Neo Energia}
    \end{itemize}
}

%\textcolor{softcol}{\hrule}

%\cvevent{2014 - 2016}{IT Consultant for IBM XPages and Notes Domino}{We4IT GmbH Bremen}{Realize projects in XPages and We4IT Aveedo, monitor project status, conduct reports}{Integrated Camunda BPMN engine and BPMN.IO modeler in We4IT Aveedo}


%\textcolor{softcol}{\hrule}

%\cvevent{2012 - 2014}{Scientific Employee / Software Development}{University of Bremen}{Invented a flexible assessment framework, targeting industrial trainees}{Supervised software development lifecycle, Recruited team members}

%\textcolor{softcol}{\hrule}

%\cvevent{2011 / 11}{Project Management Simulation Training}{Getoq Consulting}{Performed a two-day project simulation from management perspective}{Topics included customer contracts, change management, controlling, operational tasks}

%\textcolor{softcol}{\hrule}


%\cvevent{2010 - 2011}{Student Assistant / Programmer}{otulea.uni-bremen.de}{Realized an online diagnosis platform for workforce literacy development (Flex)}{Modeled software design, implemented various prototypes, conducted usability tests}


%---------------------------------------------------------------------------------------
%	EDUCATION SECTION
%--------------------------------------------------------------------------------------
\cvsection{\large{EDUCAÇÃO}}
%\cvsection{Education}

\cvevent{2023 - 2026}{\href{https://www.iesb.br/cursos/ciencia-da-computacao/}{\underline{Bacharelado em ciência da computação}}}{\href{https://www.iesb.br/home/}{\underline{Instituto de Ensino Superior de Brasília}}}{Cursando o 5º semestre, com previsão para \textbf{finalizar no 2º sem. de 2026}}{\textbf{Histórico acadêmico}: \href{https://www.icloud.com/iclouddrive/0d2XAQ77LrkM017aAy-FNSqxQ\#2025-03-11_IESB_Historico_academico_Andreas_Ramalho_Spyridakis}{\underline{\texttt{link para baixar}}}}

%\textcolor{softcol}{\hrule}

%\cvevent{2012 - 2013}{Master Project - PrIMA}{University of Bremen}{Co-Invented a touch table application for medical support, co-developed software (Java) }{Formed a scrum team, mainted project dev server (Debian), surveyed target audience}

%\textcolor{softcol}{\hrule}

%\cvevent{2012 - 2015}{Master Studies Digital Media}{University of Bremen}{Inter-cultural classes in English, covering special topics in computer science and design}{Professionalized in research methods, software development and e-assessment}

%\textcolor{softcol}{\hrule}

%\cvevent{2009 - 2010}{Semester Abroad}{University of Melbourne}{Mastered six months of study and trans-cultural experience in Melbourne, Australia}{Finished machine programming, information visualization, professional essay writing}

%---------------------------------------------------------------------------------------
%	EXTRACURRICULAR SECTION
%--------------------------------------------------------------------------------------
\cvsection{\large{EXTRACURRICULAR}}

\cvevent{03/03/2025}{\href{https://skills.yourlearning.ibm.com/activity/PLAN-3E2A749669E2}{\underline{Information Technology Fundamentals}}}{\href{https://skillsbuild.org/}{\underline{IBM SkillsBuild}}}{\textbf{Certificado de conclusão}: \href{https://www.icloud.com/iclouddrive/02bW9VQZK3lxleAA3KpnVLDew\#2025-03-06_IBM_SkillsBuild_IT_fundamentals_completion_certificate_02_Andreas}{\underline{\texttt{link para baixar}}}}{\textit{Mini-curso organizado pela IBM na plataforma SkillsBuild com o intuito de informar sobre a história da computação e conceitos essenciais de componentes de computadores, conexões de rede, hardware, software, segurança da informação e técnicas de solução de problemas (troubleshooting).}}

\cvevent{22/04/2024}{\href{https://learn.mongodb.com/courses/getting-started-with-mongodb-atlas}{\underline{Getting started with MongoDB Atlas}}}{\href{https://learn.mongodb.com/}{\underline{MongoDB University}}}{\textbf{Certificado de conclusão}: \href{https://www.icloud.com/iclouddrive/057TklO4fcJVL19epYvmqltkg\#2024-04-22_MongoDB_Andreas_Spyridakis_certificate}{\underline{\texttt{link para baixar}}}}{\textit{Curso básico para uso do banco de dados NoSQL MongoDB fornecido pela própria organização que desenvolve e mantém o banco de dados.}}

\cvevent{31/08/2023}{\href{https://www.englishtest.duolingo.com/applicants}{\underline{Teste de inglês do Duolingo}}}{\href{https://www.duolingo.com/}{\underline{Duolingo}}}{\textbf{Certificados}: \href{https://www.icloud.com/iclouddrive/070IIpYArI7B0amk6WXEumarg\#Duolingo\%5FEnglish\%5FTest-Certificado-Andreas\%5FRamalho\%5FSpyridakis}{\underline{\texttt{certificado do exame}}}}{\textbf{Resultados}: \href{https://certs.duolingo.com/7a01813affab530a8699ceb49bb3bffe}{\underline{\texttt{link para consultar}}}
}

\cvevent{Setembro de 2022}{\href{https://www.escolavirtual.gov.br/curso/629}{\underline{Aprendento com Python}}}{\href{https://www.escolavirtual.gov.br/}{\underline{Escola Virtual do Governo}}}{\textbf{Certificado de conclusão}: \href{https://www.icloud.com/iclouddrive/00fJSmWYSGsbh8dYIy4RrAQXg\#ev\%5Fg-certificado-Aprendendo\%5Fcom\%5FPython-Andreas\%5FRamalho\%5FSpyridakis}{\underline{\texttt{link para baixar}}}}{\textit{Curso rápido sobre a linguagem de programação Python com módulos que demonstram funcionalidades da linguagem como, variáveis, operações, condições, loops, funções, estruturas básicas de dados, programação orientada a objetos, usos e aplicações de bibliotecas.}}

\cvevent{Setembro de 2022}{\href{https://www.escolavirtual.gov.br/curso/74}{\underline{Sistema Eletrônico de Informações - SEI! USAR}}}{\href{https://www.escolavirtual.gov.br/}{\underline{Escola Virtual do Governo}}}{\textbf{Certificado de conclusão}: \href{https://www.icloud.com/iclouddrive/089UXLIO9A-A7_NxhjwnhKtLQ\#ev\%5Fg-certificado-Sistema\%5Feletronico\%5Fde\%5Finformacoes\%5FSEI\%5Fusar-Andreas\%5FRamalho\%5FSpyridakis}{\underline{\texttt{link para baixar}}}}{\textit{Curso rápido sobre o sistema de processos eletrônicos utilizado pelo governo do Distrito Federal e o seu uso para gerenciamento e compartilhamento de processos entre os diversos órgãos do Estado.}}

%-------------------------------------------------------------------------------------------------
%	ARTIFICIAL FOOTER (fancy footer cannot exceed linewidth) 
%--------------------------------------------------------------------------------------------------

\null
\vspace*{\fill}
\hspace{-0.25\linewidth}\colorbox{bgcol}{\makebox[1.5\linewidth][c]{\mystrut \small \textcolor{white}{andreas0r0s.carrd.co} $\cdot$ \textcolor{white}{github.com/andreas0r0s} $\cdot$ \textcolor{white}{Escrito com \LaTeX} $\cdot$ \textcolor{white}{Por Andreas R. S.}}}




%============================================================================%
%
%
%
%	DOCUMENT END
%
%
%
%============================================================================%
\end{document}
